\documentclass[11pt,a4paper]{article}
\usepackage[T1]{fontenc}
\usepackage[utf8]{inputenc}
\usepackage{lmodern,amsmath,amssymb}
\usepackage[margin=25mm]{geometry}
\usepackage[hidelinks]{hyperref}
\hypersetup{pdftitle={A shared action coefficient from classical composition},pdfauthor={navstokgap research working notes}}
\newcommand{\tightlist}{\setlength{\itemsep}{0pt}\setlength{\parskip}{0pt}}
\setlength{\emergencystretch}{3em}
\setcounter{secnumdepth}{0}
\begin{document}
\hypertarget{a-shared-action-coefficient-from-classical-composition}{%
\section{A shared action coefficient from classical
composition}\label{a-shared-action-coefficient-from-classical-composition}}

A nonnegative mass-dependent fluctuation coefficient becomes mass
independent when independent composition preserves its value as a
whole-body observable. A single finite-rate, nonzero-speed reference
constituent then makes the shared coefficient strictly positive. This
conditional result separates the composition premise from the physical
clock that sets the value.

\hypertarget{observable-and-independent-composition}{%
\subsection{1. Observable and independent
composition}\label{observable-and-independent-composition}}

All masses are positive, positions are on the line, and coefficients
have action units. For centered independent displacement variables over
a common window \(\Delta>0\), define
\(\operatorname{Var}(\Delta X_i)=\kappa_i\Delta/m_i\). Let
\(M=m_1+m_2\), \(\mu=m_1m_2/M\), \(R=(m_1X_1+m_2X_2)/M\) and
\(r=X_1-X_2\). Variance addition gives

\[\kappa_{\rm cm}:=\frac M\Delta\operatorname{Var}(\Delta R)
 =\frac{m_1\kappa_1+m_2\kappa_2}{M},\qquad
\kappa_{\rm rel}:=\frac\mu\Delta\operatorname{Var}(\Delta r)
 =\frac{m_2\kappa_1+m_1\kappa_2}{M}.\]

Thus \(\kappa_{\rm cm}+\kappa_{\rm rel}=\kappa_1+\kappa_2\) and
\(\operatorname{Cov}(\Delta R,\Delta r)=\Delta(\kappa_1-\kappa_2)/M\).
For independent Gaussian bridges replace \(\Delta\) by their common
covariance kernel; the same coefficient transformation holds. In that
Gaussian setting, centre/relative process independence is equivalent to
\(\kappa_1=\kappa_2\).

For independent stationary finite irreducible reversible chains \(J^i\),
take centered velocities \(v_i(J^i)\) and
\(X_i(t)=\int_0^t v_i(J^i_s)ds\). C019 defines
\(H_i=2m_i\int_0^\infty C_i(s)ds\). Independence gives
\(C_{\rm cm}=(m_1^2C_1+m_2^2C_2)/M^2\), so the same formulas hold with
\(H\) in place of \(\kappa\) in the long-window limit. The product state
has generator \(Q_1\otimes I+I\otimes Q_2\) and invariant law
\(\pi_1\otimes\pi_2\); detailed balance and irreducibility hold
factorwise. Retain this state even if distinct state pairs share one
velocity value: the projected velocity alone need not be Markov. If each
constituent speed is at most \(u<c\), the centre speed has the same
bound. Relative velocity may reach \(2u\); it is a coordinate
difference, not a separate particle speed premise.

\hypertarget{mass-universality-theorem}{%
\subsection{2. Mass-universality
theorem}\label{mass-universality-theorem}}

\textbf{Assumptions.} Fix a preparation class. Every mass
\(m\in(0,\infty)\) has a finite coefficient \(\kappa(m)\ge0\) depending
only on its mass within this class. Independent constituents can be
composed, and the coefficient assigned to their whole, of mass \(M\),
equals the centre coefficient calculated above. Choose either the
common-window coefficient, a Gaussian coefficient, or the Green--Kubo
plateau throughout; the theorem does not exchange these observables.

\textbf{Conclusion.} There is a constant \(K\ge0\) such that
\(\kappa(m)=K\) for every \(m>0\).

\textbf{Proof.} Set \(f(m)=m\kappa(m)\). Composition says
\(f(a+b)=f(a)+f(b)\) on positive reals. For \(b>a\), nonnegativity gives
\(f(b)-f(a)=f(b-a)\ge0\). Fix a mass unit \(m_0>0\). Additivity yields
\(f(qm_0)=qf(m_0)\) for every positive rational \(q\). For any \(m>0\),
choose rationals \(q_n\uparrow m/m_0\) and \(r_n\downarrow m/m_0\).
Monotonicity sandwiches \(f(m)\) between \(q_nf(m_0)\) and
\(r_nf(m_0)\). Hence \(f(m)=mf(m_0)/m_0\), and \(K=f(m_0)/m_0\). This
also covers \(K=0\).

The premise equating independently composed and whole-body preparations
does the physical work. All-positive-mass admissibility makes the
rational argument available; a restricted species list needs a
separately stated domain theorem.

\hypertarget{positive-reference-theorem}{%
\subsection{3. Positive reference
theorem}\label{positive-reference-theorem}}

Apply the preceding theorem to Green--Kubo plateaus. Suppose the class
contains one reference constituent of mass \(m_0\), stationary symmetric
velocities \(\pm u_0\), with \(0<u_0<c\) and reversal rate
\(0<\lambda_0<\infty\). C020 gives

\[\boxed{K=H(m_0)=\frac{m_0u_0^2}{\lambda_0}>0,
\qquad H(m)=K\quad(m>0).}\]

This excludes an attained zero in that preparation class using classical
probability, nonzero motion, a finite positive clock and the composition
premise. The value comes from the reference constituent. Other two-state
members consequently obey \(\lambda_m=mu_m^2/K\).

For any chosen \(K>0\) and common \(0<u<c\), take independent two-state
members with \(\lambda_m=mu^2/K\) and include all their finite products.
Every product has the same plateau when normalized by its total mass.
This realizes coefficient closure even though its centre velocity
generally has more than two values. It is a consistency example for the
assumptions, not a derivation of the assigned rate law. Across such
examples \(K\) can approach zero while each individual rate remains
finite. A class-wide numerical lower bound needs uniform controls, for
example \(u_0\ge u_{\min}>0\) and \(\lambda_0\le\Lambda<\infty\) at
fixed \(m_0\).

\hypertarget{preparation-and-correlation-tests}{%
\subsection{4. Preparation and correlation
tests}\label{preparation-and-correlation-tests}}

An environmental label \(\theta\) allows \(\kappa(m,\theta)=K(\theta)\)
under the same composition proof at each fixed \(\theta\). Universality
across such classes requires a preparation-equivalence premise. If cross
covariance is \(C_{12}=\operatorname{Cov}(\Delta X_1,\Delta X_2)\), the
centre coefficient has an additional \(2m_1m_2 C_{12}/(M\Delta)\);
independence fixes it to zero. Bound bodies and common reservoirs must
be tested with their actual correlations.

In A02, separately prepared tracers in otherwise identical refreshed
baths have plateau \(H(m)=(m+M_b)s^2/\nu\). Treating their total mass as
one tracer changes the coefficient by

\[H(m_a+m_b)-\frac{m_aH(m_a)+m_bH(m_b)}{m_a+m_b}
 =\frac{2m_am_bs^2}{\nu(m_a+m_b)}.\]

The positive discrepancy for \(s>0\) identifies failure of the
preparation equivalence premise, while preserving the earlier bath
calculation. The comparison is between these specified experiments;
dynamics of an actual bound composite require their own mechanical
model.

\hypertarget{scope-source-chain-and-next-test}{%
\subsection{5. Scope, source chain and next
test}\label{scope-source-chain-and-next-test}}

The theorem establishes a conditional universal positive
\textbf{long-window coefficient}, with action units and a specified
preparation class. Its role as a quantum phase constant and its
physical-time attraction are subsequent obligations. At bounded speed
the microscopic variance coefficient still tends to zero by C018. The
next A09 calculation connects these separate windows to the real
telegraph and complex checkerboard propagators.

This note develops the user's composition suggestion and the independent
P02 product-chain calculation. C019--C020 supply the reviewed
correlation input; \href{https://arxiv.org/abs/1002.4103v1}{Pavliotis
(2010)} gives its Poisson-equation source, and
\href{https://arxiv.org/abs/1802.09679v1}{Pitman--Yor (2018)} the
Gaussian covariance background. Per-result prior-art status belongs to
B16. The five dedicated algebra checks cover the two coefficients, their
sum, their cross covariance and the bath discrepancy. Existing P02
checks also test a nine-state product generator. The written proofs
supply the general additive-function and positive-reference conclusions.
\end{document}
